\subsection{Model Merging as Proxy}
We first briefly explain why model merging can be used to construct proxy models instead of training new proxy models for each sampled data mixture.

Building upon the former sections, let $\Theta_{\text{base}} \in \mathbb{R}^d$ denote the parameter vector of a pre-trained base model. Consider a collection of $N$ candidate datasets $\{D_1, D_2, \dots, D_N\}$, we define a training operator $\mathcal{T}: \mathcal{D} \times \mathbb{R}^d \to \mathbb{R}^d$, where $\mathcal{T}(D, \Theta_{\text{base}})$ results in the model parameters after training on dataset distribution $D$ initialized from $\Theta_{\text{base}}$.
For any dataset $D_i$, the \textit{parameter update vector} (or weight delta) can be defined as~\citep{yang2024model}:
\begin{equation}
    \Delta(D_i) \triangleq \mathcal{T}(D_i, \Theta_{\text{base}}) - \Theta_{\text{base}}.
\end{equation}
Consequently, the individually trained component model corresponds to parameters $\Theta_i = \Theta_{\text{base}} + \Delta(D_i)$.

The objective of data mixing is to identify an optimal mixture distribution~\citep{xie2023doremi}:
\begin{equation}
    \mathcal{D}_{\text{mix}} = \sum_{i=1}^N \alpha_i D_i, \quad \text{with } \alpha_i \ge 0 \text{ and } \sum_{i=1}^N \alpha_i = 1,
\end{equation}
so that the resulting model parameters:
\begin{equation}
    \Theta_{\text{mix}} = \mathcal{T}(\mathcal{D}_{\text{mix}}, \Theta_{\text{base}}).
\end{equation}
achieve the best performance across the target benchmarks.  


To formalize the connection between data mixing and model merging, we first state the relevant constraints.  

In practice, the magnitude of parameter updates $\Delta(D)$ remains relatively small compared to the initialization scale~\citep{gueta-etal-2023-knowledge, wu2025shadow}. Formally, for any dataset $D$ in our context, we define~\citep{wu2025shadow}:
\begin{equation}
\label{eqa:delta}
\delta=\frac{\sum|\mathcal{T}(D, \Theta_{\text{base}}) - \Theta_{\text{base}}|} {\sum|\mathcal{T}(D, \Theta_{\text{base}})| + \sum|\Theta_{\text{base}}|}\ll 1.
\end{equation}

In our experiments, $\delta$ is approximately 10\%, which satisfies this small-update assumption. Empirical studies have shown that as long as $\delta \ll 1$, the arithmetic sum of weight deltas from models trained on separate datasets closely approximates the weight delta obtained by training on their union~\citep{qin-etal-2022-exploring, wu2025shadow, lin2025efficient}:
\begin{equation}
\label{equ:satisfy}
    \Delta(D_i \cup D_j) \approx \Delta(D_i) + \Delta(D_j),
\end{equation}

% By extension, this additivity implies \textit{Homogeneity}:
% \begin{equation}
%     \Delta\left(\bigcup_{i=1}^N D_i\right) \approx \sum_{i=1}^N \Delta(D_i).
% \end{equation}
% and
% \begin{equation}
%     \Delta(\alpha D) \approx \alpha\Delta(D).
% \end{equation}

which indicates that the model parameters trained on a weighted mixture of datasets $\mathcal{D}_{\text{mix}} = \sum \alpha_i D_i$ can be approximated by the weighted average of parameters $\sum_{i=1}^N \alpha_i \Theta_i$ trained on separate datasets. We further validate the effectiveness of this approximation in our setting through an additional experiment, reported in Appendix~\ref{appd:merge_verification}.
% \begin{equation}
% \begin{aligned}
%     \Delta_{\text{mix}} &= \Delta\left(\bigcup_{i=1}^N \alpha_i D_i\right) \approx \sum_{i=1}^N \Delta(\alpha_i D_i)  \approx \sum_{i=1}^N \alpha_i \Delta(D_i).
% \end{aligned}
% \end{equation}

% Substituting the definition of the individual update vectors $\Delta(D_i) = \Theta_i - \Theta_{\text{base}}$ into the approximation:

% \begin{equation}
% \begin{aligned}
%     \Theta_{\text{mix}} &= \Theta_{\text{base}} + \Delta_{\text{mix}} \\
%     &\approx \Theta_{\text{base}} + \sum_{i=1}^N \alpha_i (\Theta_i - \Theta_{\text{base}})\\
%     &\approx \sum_{i=1}^N \alpha_i \Theta_i + \left(1 - \sum_{i=1}^N \alpha_i\right)\Theta_{\text{base}}\\
%     &\approx \sum_{i=1}^N \alpha_i \Theta_i.
% \end{aligned}
% \end{equation}

% We acknowledge that the non-linearity of training dynamics make this derivation an approximation rather than an exact identity; however, our experiments (see Table \ref{table:exp1}) demonstrates that the ranking consistency between the merged proxy and actual training remains robust, which makes the proxy sufficiently valid.

Based on this proposition, merging the component models $\{\Theta_i\}$ at a specific ratio $\{\alpha_i\}$ can obtain a proxy model for any real $\Theta_{\text{mix}}$ trained on $\{\alpha_i D_i\}$. 
% If candidate datasets are mixed with general data, to preserve equality, we will adjust the actual weights accordingly.
The merged model $M_\text{{mix}}^j$ is calculated as:
\begin{equation}
\label{merge_equation}
    M_\text{mix}^j=\sum_{i=1}^N \alpha_i^j \Theta_i,
\end{equation}
where $\Theta_i$ represents the component model of the corresponding candidate dataset $D_i$.